\IfFileExists{stacks-project.cls}{%
\documentclass{stacks-project}
}{%
\documentclass{amsart}
}

% For dealing with references we use the comment environment
\usepackage{verbatim}
\newenvironment{reference}{\comment}{\endcomment}
%\newenvironment{reference}{}{}
\newenvironment{slogan}{\comment}{\endcomment}
\newenvironment{history}{\comment}{\endcomment}

% For commutative diagrams we use Xy-pic
\usepackage[all]{xy}

% We use 2cell for 2-commutative diagrams.
\xyoption{2cell}
\UseAllTwocells

% We use multicol for the list of chapters between chapters
\usepackage{multicol}

% This is generally recommended for better output
\usepackage{lmodern}
\usepackage[T1]{fontenc}

% For cross-file-references
\usepackage{xr-hyper}

% Package for hypertext links:
\usepackage{hyperref}

% For any local file, say "hello.tex" you want to link to please
% use \externaldocument[hello-]{hello}
\externaldocument[introduction-]{introduction}
\externaldocument[conventions-]{conventions}
\externaldocument[sets-]{sets}
\externaldocument[categories-]{categories}
\externaldocument[topology-]{topology}
\externaldocument[sheaves-]{sheaves}
\externaldocument[sites-]{sites}
\externaldocument[stacks-]{stacks}
\externaldocument[fields-]{fields}
\externaldocument[algebra-]{algebra}
\externaldocument[brauer-]{brauer}
\externaldocument[homology-]{homology}
\externaldocument[derived-]{derived}
\externaldocument[simplicial-]{simplicial}
\externaldocument[more-algebra-]{more-algebra}
\externaldocument[smoothing-]{smoothing}
\externaldocument[modules-]{modules}
\externaldocument[sites-modules-]{sites-modules}
\externaldocument[injectives-]{injectives}
\externaldocument[cohomology-]{cohomology}
\externaldocument[sites-cohomology-]{sites-cohomology}
\externaldocument[dga-]{dga}
\externaldocument[dpa-]{dpa}
\externaldocument[sdga-]{sdga}
\externaldocument[hypercovering-]{hypercovering}
\externaldocument[schemes-]{schemes}
\externaldocument[constructions-]{constructions}
\externaldocument[properties-]{properties}
\externaldocument[morphisms-]{morphisms}
\externaldocument[coherent-]{coherent}
\externaldocument[divisors-]{divisors}
\externaldocument[limits-]{limits}
\externaldocument[varieties-]{varieties}
\externaldocument[topologies-]{topologies}
\externaldocument[descent-]{descent}
\externaldocument[perfect-]{perfect}
\externaldocument[more-morphisms-]{more-morphisms}
\externaldocument[flat-]{flat}
\externaldocument[groupoids-]{groupoids}
\externaldocument[more-groupoids-]{more-groupoids}
\externaldocument[etale-]{etale}
\externaldocument[chow-]{chow}
\externaldocument[intersection-]{intersection}
\externaldocument[pic-]{pic}
\externaldocument[weil-]{weil}
\externaldocument[adequate-]{adequate}
\externaldocument[dualizing-]{dualizing}
\externaldocument[duality-]{duality}
\externaldocument[discriminant-]{discriminant}
\externaldocument[derham-]{derham}
\externaldocument[local-cohomology-]{local-cohomology}
\externaldocument[algebraization-]{algebraization}
\externaldocument[curves-]{curves}
\externaldocument[resolve-]{resolve}
\externaldocument[models-]{models}
\externaldocument[functors-]{functors}
\externaldocument[equiv-]{equiv}
\externaldocument[pione-]{pione}
\externaldocument[etale-cohomology-]{etale-cohomology}
\externaldocument[proetale-]{proetale}
\externaldocument[relative-cycles-]{relative-cycles}
\externaldocument[more-etale-]{more-etale}
\externaldocument[trace-]{trace}
\externaldocument[crystalline-]{crystalline}
\externaldocument[spaces-]{spaces}
\externaldocument[spaces-properties-]{spaces-properties}
\externaldocument[spaces-morphisms-]{spaces-morphisms}
\externaldocument[decent-spaces-]{decent-spaces}
\externaldocument[spaces-cohomology-]{spaces-cohomology}
\externaldocument[spaces-limits-]{spaces-limits}
\externaldocument[spaces-divisors-]{spaces-divisors}
\externaldocument[spaces-over-fields-]{spaces-over-fields}
\externaldocument[spaces-topologies-]{spaces-topologies}
\externaldocument[spaces-descent-]{spaces-descent}
\externaldocument[spaces-perfect-]{spaces-perfect}
\externaldocument[spaces-more-morphisms-]{spaces-more-morphisms}
\externaldocument[spaces-flat-]{spaces-flat}
\externaldocument[spaces-groupoids-]{spaces-groupoids}
\externaldocument[spaces-more-groupoids-]{spaces-more-groupoids}
\externaldocument[bootstrap-]{bootstrap}
\externaldocument[spaces-pushouts-]{spaces-pushouts}
\externaldocument[spaces-chow-]{spaces-chow}
\externaldocument[groupoids-quotients-]{groupoids-quotients}
\externaldocument[spaces-more-cohomology-]{spaces-more-cohomology}
\externaldocument[spaces-simplicial-]{spaces-simplicial}
\externaldocument[spaces-duality-]{spaces-duality}
\externaldocument[formal-spaces-]{formal-spaces}
\externaldocument[restricted-]{restricted}
\externaldocument[spaces-resolve-]{spaces-resolve}
\externaldocument[formal-defos-]{formal-defos}
\externaldocument[defos-]{defos}
\externaldocument[cotangent-]{cotangent}
\externaldocument[examples-defos-]{examples-defos}
\externaldocument[algebraic-]{algebraic}
\externaldocument[examples-stacks-]{examples-stacks}
\externaldocument[stacks-sheaves-]{stacks-sheaves}
\externaldocument[criteria-]{criteria}
\externaldocument[artin-]{artin}
\externaldocument[quot-]{quot}
\externaldocument[stacks-properties-]{stacks-properties}
\externaldocument[stacks-morphisms-]{stacks-morphisms}
\externaldocument[stacks-limits-]{stacks-limits}
\externaldocument[stacks-cohomology-]{stacks-cohomology}
\externaldocument[stacks-perfect-]{stacks-perfect}
\externaldocument[stacks-introduction-]{stacks-introduction}
\externaldocument[stacks-more-morphisms-]{stacks-more-morphisms}
\externaldocument[stacks-geometry-]{stacks-geometry}
\externaldocument[moduli-]{moduli}
\externaldocument[moduli-curves-]{moduli-curves}
\externaldocument[examples-]{examples}
\externaldocument[exercises-]{exercises}
\externaldocument[guide-]{guide}
\externaldocument[desirables-]{desirables}
\externaldocument[coding-]{coding}
\externaldocument[obsolete-]{obsolete}
\externaldocument[fdl-]{fdl}
\externaldocument[index-]{index}

% Theorem environments.
%
\theoremstyle{plain}
\newtheorem{theorem}[subsection]{Theorem}
\newtheorem{proposition}[subsection]{Proposition}
\newtheorem{lemma}[subsection]{Lemma}

\theoremstyle{definition}
\newtheorem{definition}[subsection]{Definition}
\newtheorem{example}[subsection]{Example}
\newtheorem{exercise}[subsection]{Exercise}
\newtheorem{situation}[subsection]{Situation}

\theoremstyle{remark}
\newtheorem{remark}[subsection]{Remark}
\newtheorem{remarks}[subsection]{Remarks}

\numberwithin{equation}{subsection}

% Macros
%
\def\lim{\mathop{\mathrm{lim}}\nolimits}
\def\colim{\mathop{\mathrm{colim}}\nolimits}
\def\Spec{\mathop{\mathrm{Spec}}}
\def\Hom{\mathop{\mathrm{Hom}}\nolimits}
\def\Ext{\mathop{\mathrm{Ext}}\nolimits}
\def\SheafHom{\mathop{\mathcal{H}\!\mathit{om}}\nolimits}
\def\SheafExt{\mathop{\mathcal{E}\!\mathit{xt}}\nolimits}
\def\Sch{\mathit{Sch}}
\def\Mor{\mathop{\mathrm{Mor}}\nolimits}
\def\Ob{\mathop{\mathrm{Ob}}\nolimits}
\def\Sh{\mathop{\mathit{Sh}}\nolimits}
\def\NL{\mathop{N\!L}\nolimits}
\def\CH{\mathop{\mathrm{CH}}\nolimits}
\def\proetale{{pro\text{-}\acute{e}tale}}
\def\etale{{\acute{e}tale}}
\def\QCoh{\mathit{QCoh}}
\def\Ker{\mathop{\mathrm{Ker}}}
\def\Im{\mathop{\mathrm{Im}}}
\def\Coker{\mathop{\mathrm{Coker}}}
\def\Coim{\mathop{\mathrm{Coim}}}

% Boxtimes
%
\DeclareMathSymbol{\boxtimes}{\mathbin}{AMSa}{"02}

%
% Macros for moduli stacks/spaces
%
\def\QCohstack{\mathcal{QC}\!\mathit{oh}}
\def\Cohstack{\mathcal{C}\!\mathit{oh}}
\def\Spacesstack{\mathcal{S}\!\mathit{paces}}
\def\Quotfunctor{\mathrm{Quot}}
\def\Hilbfunctor{\mathrm{Hilb}}
\def\Curvesstack{\mathcal{C}\!\mathit{urves}}
\def\Polarizedstack{\mathcal{P}\!\mathit{olarized}}
\def\Complexesstack{\mathcal{C}\!\mathit{omplexes}}
% \Pic is the operator that assigns to X its picard group, usage \Pic(X)
% \Picardstack_{X/B} denotes the Picard stack of X over B
% \Picardfunctor_{X/B} denotes the Picard functor of X over B
\def\Pic{\mathop{\mathrm{Pic}}\nolimits}
\def\Picardstack{\mathcal{P}\!\mathit{ic}}
\def\Picardfunctor{\mathrm{Pic}}
\def\Deformationcategory{\mathcal{D}\!\mathit{ef}}
\usepackage{cleveref}
\crefname{originallemma}{original lemma}{original lemmas}
\Crefname{originallemma}{Original lemma}{Original lemmas}
\makeatletter
\newlabel{more-morphisms-lemma-deform-projective}{{0D4F}{source}{}{sourcepointer-1}{}}
\newlabel{more-morphisms-lemma-deform-projective@cref}{{[originallemma][0][]0D4F}{[1][0][]source}{}{}{}}
\makeatother
\crefname{originallemma}{original lemma}{original lemmas}
\Crefname{originallemma}{Original lemma}{Original lemmas}
\makeatletter
\newlabel{more-algebra-lemma-lim-finite-projective-gives-finite-projective}{{0D4B}{source}{}{sourcepointer-2}{}}
\newlabel{more-algebra-lemma-lim-finite-projective-gives-finite-projective@cref}{{[originallemma][0][]0D4B}{[1][0][]source}{}{}{}}
\makeatother
\crefname{originallemma}{original lemma}{original lemmas}
\Crefname{originallemma}{Original lemma}{Original lemmas}
\makeatletter
\newlabel{divisors-lemma-relative-proj-flat}{{0D4C}{source}{}{sourcepointer-3}{}}
\newlabel{divisors-lemma-relative-proj-flat@cref}{{[originallemma][0][]0D4C}{[1][0][]source}{}{}{}}
\makeatother
\crefname{originallemma}{original lemma}{original lemmas}
\Crefname{originallemma}{Original lemma}{Original lemmas}
\makeatletter
\newlabel{constructions-lemma-base-change-map-proj}{{01N2}{source}{}{sourcepointer-4}{}}
\newlabel{constructions-lemma-base-change-map-proj@cref}{{[originallemma][0][]01N2}{[1][0][]source}{}{}{}}
\makeatother
\crefname{originallemma}{original lemma}{original lemmas}
\Crefname{originallemma}{Original lemma}{Original lemmas}
\makeatletter
\newlabel{divisors-lemma-relative-proj-proper}{{0800}{source}{}{sourcepointer-5}{}}
\newlabel{divisors-lemma-relative-proj-proper@cref}{{[originallemma][0][]0800}{[1][0][]source}{}{}{}}
\makeatother
\crefname{originallemma}{original lemma}{original lemmas}
\Crefname{originallemma}{Original lemma}{Original lemmas}
\makeatletter
\newlabel{divisors-lemma-relative-proj-finite-presentation}{{0D4D}{source}{}{sourcepointer-6}{}}
\newlabel{divisors-lemma-relative-proj-finite-presentation@cref}{{[originallemma][0][]0D4D}{[1][0][]source}{}{}{}}
\makeatother
\crefname{originallemma}{original lemma}{original lemmas}
\Crefname{originallemma}{Original lemma}{Original lemmas}
\makeatletter
\newlabel{more-morphisms-lemma-finite-free-open}{{0CZR}{source}{}{sourcepointer-7}{}}
\newlabel{more-morphisms-lemma-finite-free-open@cref}{{[originallemma][0][]0CZR}{[1][0][]source}{}{}{}}
\makeatother
\crefname{originallemma}{original lemma}{original lemmas}
\Crefname{originallemma}{Original lemma}{Original lemmas}
\makeatletter
\newlabel{more-algebra-lemma-limit-henselian}{{0CT7}{source}{}{sourcepointer-8}{}}
\newlabel{more-algebra-lemma-limit-henselian@cref}{{[originallemma][0][]0CT7}{[1][0][]source}{}{}{}}
\makeatother
\crefname{originallemma}{original lemma}{original lemmas}
\Crefname{originallemma}{Original lemma}{Original lemmas}
\makeatletter
\newlabel{more-morphisms-lemma-ample-in-neighbourhood}{{0D2S}{source}{}{sourcepointer-9}{}}
\newlabel{more-morphisms-lemma-ample-in-neighbourhood@cref}{{[originallemma][0][]0D2S}{[1][0][]source}{}{}{}}
\makeatother
\crefname{originallemma}{original lemma}{original lemmas}
\Crefname{originallemma}{Original lemma}{Original lemmas}
\makeatletter
\newlabel{algebra-lemma-NAK}{{00DV}{source}{}{sourcepointer-10}{}}
\newlabel{algebra-lemma-NAK@cref}{{[originallemma][0][]00DV}{[1][0][]source}{}{}{}}
\makeatother
\title{Strong formal effectiveness for polarized proper flat finitely presented schemes}
\author{The Stacks Project Authors}
\date{stacks-project:quot.tex:local-snapshot-20261012; GNU FDL 1.2 or later}
\begin{document}
\maketitle
\setcounter{subsection}{1}
\noindent\textit{Editorial note.} All delivered mathematical text is the author's own, copied from the pinned Stacks Project chapter file quot.tex; no statement or proof step was rewritten, reordered or removed. The counted span is the author environment and its own proof environment. References inside the span are handed to the assembler as navigation stubs that carry the canonical tag of the cited tag only; whether each of them is load-bearing is recorded on the stub itself.\par\medskip
\noindent\textit{Allowed standard prerequisites (editorial).}\begin{itemize}
\item Lemma more-morphisms-lemma-deform-projective (canonical tag 0D4F). cited at proof line 4544. Supplies d\_0, finite presentation of A\_n, projective homogeneous parts, and system compatibility.
\item Lemma more-algebra-lemma-lim-finite-projective-gives-finite-projective (canonical tag 0D4B). cited at proof lines 4565-4566. B\_d finite projective and B x\_R R\_n = A\_n.
\item Lemma divisors-lemma-relative-proj-flat (canonical tag 0D4C). cited at proof lines 4575-4576. Proj(B) flat over S, L flat.
\item Lemma constructions-lemma-base-change-map-proj (canonical tag 01N2). cited at proof lines 4577-4578. Proj commutes with base change, giving the descent isomorphisms.
\item Lemma divisors-lemma-relative-proj-proper (canonical tag 0800). cited at proof lines 4588-4589. Properness of X -> S.
\item Lemma divisors-lemma-relative-proj-finite-presentation (canonical tag 0D4D). cited at proof lines 4589-4590. Finite presentation of X -> S and of L.
\item Lemma more-morphisms-lemma-finite-free-open (canonical tag 0CZR). cited at proof lines 4594-4595. Invertibility locus is an open neighbourhood of X\_1.
\item Lemma more-algebra-lemma-limit-henselian (canonical tag 0CT7). cited at proof lines 4598-4599 and 4619-4620. Ker(R -> R\_1) in the Jacobson radical.
\item Lemma more-morphisms-lemma-ample-in-neighbourhood (canonical tag 0D2S). cited at proof lines 4602-4603. Ampleness on fibres is open and contains S\_1.
\item Lemma algebra-lemma-NAK (canonical tag 00DV). cited at proof line 4618, used again at 4623-4624. Nakayama: Psi surjective and Ker(Psi) generated by the lifted relations.
\end{itemize}\medskip
\ifdefined\hypertarget\hypertarget{sourcepointer-1}{}\fi
\noindent\textit{Source reference.} Original lemma 0D4F: sources/stacks-chow/tags.txt, source lines 16450--16450. Cited at proof line 4544: gives d\_0 and the finite-presentation properties of each A\_n plus compatibility of the system; load-bearing.\par\medskip
\ifdefined\hypertarget\hypertarget{sourcepointer-2}{}\fi
\noindent\textit{Source reference.} Original lemma 0D4B: sources/stacks-chow/tags.txt, source lines 16446--16446. Cited at proof lines 4565-4566: B\_d finite projective over R and B x\_R R\_n = A\_n; load-bearing.\par\medskip
\ifdefined\hypertarget\hypertarget{sourcepointer-3}{}\fi
\noindent\textit{Source reference.} Original lemma 0D4C: sources/stacks-chow/tags.txt, source lines 16447--16447. Cited at proof lines 4575-4576: X = Proj(B) flat over S and L flat; load-bearing.\par\medskip
\ifdefined\hypertarget\hypertarget{sourcepointer-4}{}\fi
\noindent\textit{Source reference.} Original lemma 01N2: sources/stacks-chow/tags.txt, source lines 2143--2143. Cited at proof lines 4577-4578: canonical isomorphisms X x\_S S\_n = X\_n and L|\_\{X\_n\} = L\_n; load-bearing.\par\medskip
\ifdefined\hypertarget\hypertarget{sourcepointer-5}{}\fi
\noindent\textit{Source reference.} Original lemma 0800: sources/stacks-chow/tags.txt, source lines 10165--10165. Cited at proof lines 4588-4589: B finite presentation over R makes X -> S proper; load-bearing.\par\medskip
\ifdefined\hypertarget\hypertarget{sourcepointer-6}{}\fi
\noindent\textit{Source reference.} Original lemma 0D4D: sources/stacks-chow/tags.txt, source lines 16448--16448. Cited at proof lines 4589-4590: X -> S of finite presentation and L finitely presented; load-bearing.\par\medskip
\ifdefined\hypertarget\hypertarget{sourcepointer-7}{}\fi
\noindent\textit{Source reference.} Original lemma 0CZR: sources/stacks-chow/tags.txt, source lines 16286--16286. Cited at proof lines 4594-4595: L is invertible on an open neighbourhood of X\_1; load-bearing.\par\medskip
\ifdefined\hypertarget\hypertarget{sourcepointer-8}{}\fi
\noindent\textit{Source reference.} Original lemma 0CT7: sources/stacks-chow/tags.txt, source lines 16057--16057. Cited at proof lines 4598-4599 and 4619-4620: Ker(R -> R\_1) inside the Jacobson radical, making every neighbourhood of X\_1 equal to X and enabling Nakayama; load-bearing twice.\par\medskip
\ifdefined\hypertarget\hypertarget{sourcepointer-9}{}\fi
\noindent\textit{Source reference.} Original lemma 0D2S: sources/stacks-chow/tags.txt, source lines 16392--16392. Cited at proof lines 4602-4603: ampleness locus open and contains S\_1, hence equals S; load-bearing.\par\medskip
\ifdefined\hypertarget\hypertarget{sourcepointer-10}{}\fi
\noindent\textit{Source reference.} Original lemma 00DV: sources/stacks-chow/tags.txt, source lines 512--512. Cited at proof line 4618 (and used again at 4623-4624): surjectivity of Psi and generation of its kernel by the lifts G\_j; load-bearing twice.\par\medskip
\begin{lemma}[Strong formal effectiveness for polarized schemes]
\label{lemma-polarized-strong-effectiveness}
\begin{slogan}
Grothendieck's algebraization theorem continues to hold in
the non-Noetherian setting if one assumes flatness and
finite presentation.
\end{slogan}
Let $(R_n)$ be an inverse system of rings with surjective transition maps
whose kernels are locally nilpotent. Set $R = \lim R_n$.
Set $S_n = \Spec(R_n)$ and $S = \Spec(R)$. Consider a commutative diagram
$$
\xymatrix{
X_1 \ar[r]_{i_1} \ar[d] & X_2 \ar[r]_{i_2} \ar[d] & X_3 \ar[r] \ar[d] &
\ldots \\
S_1 \ar[r] & S_2 \ar[r] & S_3 \ar[r] & \ldots
}
$$
of schemes with cartesian squares. Suppose given $(\mathcal{L}_n, \varphi_n)$
where each $\mathcal{L}_n$ is an invertible sheaf on $X_n$ and
$\varphi_n : i_n^*\mathcal{L}_{n + 1} \to \mathcal{L}_n$ is an isomorphism.
If
\begin{enumerate}
\item $X_n \to S_n$ is proper, flat, of finite presentation, and
\item $\mathcal{L}_1$ is ample on $X_1$
\end{enumerate}
then there exists a morphism of schemes $X \to S$
proper, flat, and of finite presentation
and an ample invertible $\mathcal{O}_X$-module $\mathcal{L}$
and isomorphisms $X_n \cong X \times_S S_n$ and
$\mathcal{L}_n \cong \mathcal{L}|_{X_n}$ compatible with
the morphisms $i_n$ and $\varphi_n$.
\end{lemma}

\begin{proof}
Choose $d_0$ for $X_1 \to S_1$ and $\mathcal{L}_1$ as in
More on Morphisms, Lemma \ref{more-morphisms-lemma-deform-projective}.
For any $n \geq 1$ set
$$
A_n = R_n \oplus
\bigoplus\nolimits_{d \geq d_0} H^0(X_n, \mathcal{L}_n^{\otimes d})
$$
By the lemma each $A_n$ is a finitely presented graded $R_n$-algebra
whose homogeneous parts $(A_n)_d$ are finite projective $R_n$-modules
such that $X_n = \text{Proj}(A_n)$ and
$\mathcal{L}_n = \mathcal{O}_{\text{Proj}(A_n)}(1)$.
The lemma also guarantees that the maps
$$
A_1 \leftarrow A_2 \leftarrow A_3 \leftarrow \ldots
$$
induce isomorphisms $A_n = A_m \otimes_{R_m} R_n$ for $n \leq m$.
We set
$$
B = \bigoplus\nolimits_{d \geq 0} B_d
\quad\text{with}\quad
B_d = \lim_n (A_n)_d
$$
By More on Algebra, Lemma
\ref{more-algebra-lemma-lim-finite-projective-gives-finite-projective}
we see that $B_d$ is a finite projective $R$-module and that
$B \otimes_R R_n = A_n$. Thus the scheme
$$
X = \text{Proj}(B)
\quad\text{and}\quad
\mathcal{L} = \mathcal{O}_X(1)
$$
is flat over $S$ and $\mathcal{L}$ is a quasi-coherent $\mathcal{O}_X$-module
flat over $S$, see
Divisors, Lemma \ref{divisors-lemma-relative-proj-flat}.
Because formation of Proj commutes with base change
(Constructions, Lemma \ref{constructions-lemma-base-change-map-proj})
we obtain canonical isomorphisms
$$
X \times_S S_n = X_n
\quad\text{and}\quad
\mathcal{L}|_{X_n} \cong \mathcal{L}_n
$$
compatible with the transition maps of the system.
Thus we may think of $X_1 \subset X$ as a closed subscheme.
Below we will show that $B$ is of finite presentation over $R$.
By Divisors, Lemmas \ref{divisors-lemma-relative-proj-proper} and
\ref{divisors-lemma-relative-proj-finite-presentation}
this implies that $X \to S$ is of finite presentation
and proper and that $\mathcal{L} = \mathcal{O}_X(1)$
is of finite presentation as an $\mathcal{O}_X$-module.
Since the restriction of $\mathcal{L}$ to the base change
$X_1 \to S_1$ is invertible, we see from
More on Morphisms, Lemma \ref{more-morphisms-lemma-finite-free-open}
that $\mathcal{L}$ is invertible on an open neighbourhood of $X_1$ in $X$.
Since $X \to S$ is closed and since $\Ker(R \to R_1)$
is contained in the Jacobson radical
(More on Algebra, Lemma \ref{more-algebra-lemma-limit-henselian})
we see that any open neighbourhood of $X_1$ in $X$ is equal to $X$.
Thus $\mathcal{L}$ is invertible. Finally, the set of points in
$S$ where $\mathcal{L}$ is ample on the fibre is open in $S$
(More on Morphisms, Lemma \ref{more-morphisms-lemma-ample-in-neighbourhood})
and contains $S_1$ hence equals $S$. Thus $X \to S$ and $\mathcal{L}$
have all the properties required of them in the statement of the lemma.

\medskip\noindent
We prove the claim above.
Choose a presentation $A_1 = R_1[X_1, \ldots, X_s]/(F_1, \ldots, F_t)$
where $X_i$ are variables having degrees $d_i$ and $F_j$
are homogeneous polynomials in $X_i$ of degree $e_j$.
Then we can choose a map
$$
\Psi : R[X_1, \ldots, X_s] \longrightarrow B
$$
lifting the map $R_1[X_1, \ldots, X_s] \to A_1$. Since each $B_d$
is finite projective over $R$ we conclude from 
Nakayama's lemma (Algebra, Lemma \ref{algebra-lemma-NAK}
using again that $\Ker(R \to R_1)$ is contained in the Jacobson radical
of $R$) that $\Psi$ is surjective. Since $- \otimes_R R_1$ is right
exact we can find $G_1, \ldots, G_t \in \Ker(\Psi)$
mapping to $F_1, \ldots, F_t$ in $R_1[X_1, \ldots, X_s]$.
Observe that $\Ker(\Psi)_d$ is a finite projective $R$-module
for all $d \geq 0$ as the kernel of the surjection
$R[X_1, \ldots, X_s]_d \to B_d$ of finite projective $R$-modules.
We conclude from Nakayama's lemma once more that 
$\Ker(\Psi)$ is generated by $G_1, \ldots, G_t$.
\end{proof}

\end{document}
